% Deterministically prepared from the preserved v107 source; see prepare.py.
\section{Endpoint fans, reconstruction lemmas, and general fibres}\label{sec:endpoint-completion}
The forward active-set formulas are classical quadratic programming. The issue here is the exhaustion of inverse representations when the active labels are not given. We expose the two-dimensional fan and then separate the reconstruction arguments used on the first nonvisible strata.

\begin{lemma}[Planar complementary fans]\label{lem:planar-fan}
Let $Y\subset\R^2$ be an open convex cone obtained as the image of $X^\circ$ under two independent linear forms $U,V$. In the chain normalization
\[
 B=[u,au+v],\qquad C=\begin{pmatrix}1&a\\a&1+a^2\end{pmatrix},
\]
the four strict primal active sectors are exactly
\[
\begin{array}{c|c|c}
 I&\text{strict inequalities}&(z_1,z_2)\\\hline
 \varnothing&U>0,\ aU+V>0&(0,0)\\
 \{1\}&U<0,\ V>0&(-U,0)\\
 \{2\}&aU+V<0,\ U-aV>0&(0,-(aU+V)/(1+a^2))\\
 \{1,2\}&V<0,\ U-aV<0&(-U+aV,-V)
\end{array}
\]
If the empty, first-only, and full sectors meet $Y$ and the second-only sector does not, then $Y$ contains no point with $U>0,V<0$. On the retained cone the value function is exactly the chain form \eqref{eq:chain-normal}; feasibility is equivalent to the two ratio bounds in \eqref{eq:chain-fibre}.

For the dual maximization $\max_{w\ge0}\{2(U,V)w-w^{\mathsf T}D(a)w\}$, where $D(a)=\left(\begin{smallmatrix}1&a\\a&1\end{smallmatrix}\right)$ and $|a|<1$, the strict sectors are
\[
\begin{array}{c|c|c}
 J&\text{strict inequalities}&(w_1,w_2)\\\hline
 \varnothing&U<0,\ V<0&(0,0)\\
 \{1\}&U>0,\ V-aU<0&(U,0)\\
 \{2\}&V>0,\ U-aV<0&(0,V)\\
 \{1,2\}&U-aV>0,\ V-aU>0&D(a)^{-1}(U,V)^{\mathsf T}
\end{array}
\]
If the first three sectors meet $Y$ and the last does not, the regions $U>0$ and $V>0$ are disjoint in $Y$. The value is then $U_+^2+V_+^2$, with exactly the bound $a\ge a_0$ in \eqref{eq:fork-fibre}. Replacing $(U,V)$ by $(-U,-V)$ gives the missing-full primal case.
\end{lemma}
\begin{proof}
Each row follows by solving stationarity on the indicated positive coordinates and requiring the other multipliers to be positive. These equations prove both necessity and sufficiency of every row, including uniqueness by strict convexity. Weak versions cover the sector boundaries.

For the chain assertion, suppose first that $a\ge0$ and $Y$ contains $U>0,V<0$. At such a point $U-aV>0$. Avoidance of the missing sector forces $aU+V\ge0$. Equality is impossible in the open set $Y$, since a small perturbation enters the missing sector. When $a=0$ the strict inequality is already impossible. When $a>0$, join this point to a strict full-sector point. The second coordinate remains negative along the segment; at the full-sector endpoint both $U$ and $aU+V$ are negative. At the crossing $aU+V=0$ one has $U=-V/a>0$, so immediately beyond it the segment lies in the missing sector, a contradiction.

If $a<0$, a point with $U>0,V<0$ has $aU+V<0$. Avoidance of the missing sector forces $U-aV<0$, again strictly by openness. Join it to a strict empty-sector point. The first coordinate stays positive. At the crossing $U-aV=0$ one has $V=U/a<0$ and $aU+V<0$; the segment therefore crosses the missing sector. This is again impossible.

Thus $V<0$ implies $U<0$ on $Y$. The three candidate optimizers in the completed-square coordinates $w_1=z_1+az_2,w_2=z_2$ are $(0,0)$, $(-U,0)$, and $(-U,-V)$. For $U>0$ the remaining inactive inequality is $aU+V\ge0$, equivalent to $a\ge -V/U$. For $V<0$ the remaining feasibility inequality is $-U+aV\ge0$, equivalent to $a\le U/V$. For $U<0<V$, stationarity and the inactive multiplier are already satisfied. These are all possibilities off the two separating lines. They prove necessity and sufficiency of the displayed supremum and infimum bounds. Continuity extends the equality to the boundary of $X$.

For the dual assertion, if $a\le0$ the positive quadrant lies in the both-active sector. Suppose $0<a<1$ and a point of $Y$ has $U,V>0$ but lies outside that sector. It is in one singleton sector, say $V<aU$. Join it to a strict point of the other singleton sector, which has $V>0$ and $U<aV$. The second coordinate stays positive and the continuous ratio $U/V$ crosses the interval $(a,1/a)$. On this interval both-active inequalities hold. This contradicts the missing-sector hypothesis. Exchanging coordinates treats the other possibility. The only optimizers on $Y$ are therefore $(0,0),(U,0),(0,V)$. Their remaining multiplier inequalities are precisely $V\le aU$ on $U>0$ and $U\le aV$ on $V>0$. These give the stated ratio bound and show its sufficiency. Sign reversal gives the primal missing-full assertion.
\end{proof}

\begin{proposition}[Exhaustion of the three-cell inverse fibre]\label{prop:planar-exhaustion}
Under the hypotheses of Theorem~\ref{thm:three-cells}, every minimal competitor is one of the three families displayed there. The cell labels determine the signs of the two rank-one forms, the chain or fork order type, and the corresponding global gluing on $X$. There is no further gluing choice with the same labelled function.
\end{proposition}
\begin{proof}
Column independence gives four distinct algebraic active-set matrices, ordered as the Boolean lattice on two vertices by Lemma~\ref{lem:lattice}. Removing a singleton leaves a three-element chain; removing the empty or full vertex leaves one of the two forks. In each case the ranges of two rank-one adjacent differences span a two-dimensional space. Every representation of the same function has a cross block containing those ranges, so it needs at least two endpoints. A two-endpoint competitor therefore has independent cross columns. It has exactly three visible matrices, by the strict-point argument of Lemma~\ref{lem:competitor}; hence it has the same chain or fork order type.

For a chain, normalize its first diagonal pivot and its second Schur pivot to one by positive endpoint scaling. Completing the square then forces $C=\left(\begin{smallmatrix}1&a\\a&1+a^2\end{smallmatrix}\right)$ and $B=[u,au+v]$. The two successive rank-one deficits determine $uu^{\mathsf T}$ and $vv^{\mathsf T}$. Their orientations are fixed by the side on which the respective decrement occurs; reversing either sign would put a different polynomial on a labelled open set. Lemma~\ref{lem:planar-fan} proves exactly the global nesting and interval bounds. The full Schur complement is the observed smallest matrix, so positivity adds no hidden condition.

For a missing-empty fork use $S=S_{\{1,2\}}$, $D=C^{-1}$ and $H=BC^{-1}$. Completing the square and dualizing gives
\[
 G(x)=x^{\mathsf T}Sx+\max_{w\ge0}\{2(H^{\mathsf T}x)^{\mathsf T}w-w^{\mathsf T}Dw\}.
\]
Normalize the diagonal of $D$ to one. The two observed increments above $S$ determine $u,v$, positively oriented on their increment regions. The dual table forces exactly $D(a)$ and its ratio bound, and reversing the completion of the square gives \eqref{eq:fork-fibre}. Its Schur complement is $S>0$. The missing-full case is the sign-reversed primal table; its unobserved full Schur complement must be positive, which is exactly the additional inequality in Theorem~\ref{thm:three-cells}.

All sector conditions were derived from the unique optimizer and were shown sufficient on the whole cone, not only near a witness point. Every minimal competitor must therefore have the same oriented rank-one labels and one of the stated scalar parameters. This proves exhaustion and excludes alternative gluing. On an admissible interval reduced to a single point, only the two endpoint scales remain; on a nontrivial open interval they coexist with the scalar modulus, giving dimension three. Endpoint permutations remain a discrete ambiguity.
\end{proof}

\begin{lemma}[Reconstruction of a missing singleton]\label{lem:missing-singleton}
Let $e\ge3$, $\rank B=e$, and suppose the only missing active set is a singleton. The observed data determine the normalized missing cross column and then the entire normalized representation.
\end{lemma}
\begin{proof}
The largest and smallest matrices $A,S$ are present. Exactly $e-1$ rank-one covers of $A$ remain; factor their deficits as $u_iu_i^{\mathsf T}$ and orient $u_i$ positively in the empty-active region. A rank-two deficit whose range contains exactly one of the known column lines must correspond to the pair consisting of that column and the missing column. Indeed the range of a deficit is the span of exactly its active columns, and independent columns have no additional known column line in their span. Two such pair ranges intersect in the missing line, since their three original columns are independent. Such two pairs exist because $e\ge3$.

Choose a unit vector $v$ on that line with positive orientation on the empty region. For a recovered pair express its deficit uniquely as $[u_i,v]K[u_i,v]^{\mathsf T}$. If the missing normalized column is $d v$, its normalized two-endpoint correlation is some $a$ with $|a|<1$, and
\[
 K=\frac1{1-a^2}\begin{pmatrix}1&-ad\\-ad&d^2\end{pmatrix}.
\]
Thus $d=\sqrt{K_{22}/K_{11}}$ is intrinsic and positive. All normalized columns are now known; $A-S$ and formula \eqref{eq:endpoint-reconstruct} determine the coupling. This construction uses only observed labels and ranges, so it is forced for every minimal competitor.
\end{proof}

\begin{lemma}[Reconstruction of a missing extreme]\label{lem:missing-extreme}
Let $e\ge3$ and $\rank B=e$. If exactly the full active cell is missing, the observed singleton and pair deficits below $A$ determine the full normalized representation. If exactly the empty cell is missing, the dual singleton and pair increments above $S_{[e]}$ do so.
\end{lemma}
\begin{proof}
For the first assertion, singleton deficits determine the positively oriented normalized columns $u_j=B_j/\sqrt{C_{jj}}$. Every pair is observed since $e\ge3$. In the pair basis $U_{ij}=[u_i,u_j]$, its deficit has coefficient matrix $\overline C_{\{i,j\},\{i,j\}}^{-1}$. A left inverse of $U_{ij}$ recovers this coefficient, and inversion gives the corresponding normalized principal block. All pairs determine every off-diagonal entry and the diagonal is one. Together with the observed $A$ this gives the whole matrix.

For the second assertion write $S=S_{[e]}$, $D=C^{-1}$, and $H=BC^{-1}$. Block elimination gives
\[
 S_I-S=H_JD_{JJ}^{-1}H_J^{\mathsf T},\qquad J=I^c.
\]
Every dual singleton and pair is present. Normalize $h_j=H_j/\sqrt{D_{jj}}$ and orient it negatively on the full-active region, where $C^{-1}B^{\mathsf T}x<0$. The preceding pair argument recovers $\overline D$ with diagonal one. Then
\[
 C_*=\overline D^{-1},\quad B_*=(h_j)C_*,\quad
 A_*=S+(h_j)C_*(h_j)^{\mathsf T}
\]
is the unique normalized representative, with $\diag C_*^{-1}=\mathbf1$. Multiple maximal or minimal observed vertices distinguish the missing extremes intrinsically, so the construction does not require supplied active-set names.
\end{proof}

Together with the unchanged-extremes reconstruction when the missing vertex is neither an extreme nor a singleton, these lemmas cover every case of Theorem~\ref{thm:one-missing}. The rank of the span of observed Hessian differences is $e$, even when an extreme is absent, because the ranges adjacent to that extreme span all cross columns. This supplies the common minimality and competitor argument before any reconstruction is performed.

\subsection{A finite presentation beyond the first strata}
An explicit scalar normal form is more informative than a quantifier-elimination output. Nevertheless arbitrary simultaneous cell disappearance and rank loss admit the following finite inverse description. It treats all active cones, including those with empty interior.

\begin{theorem}[Semialgebraic presentation of all minimal endpoint fibres]\label{thm:all-strata}
This is a reference-presentation equality theorem: the matrix $Q$ is supplied. The function-only minimal-fibre theorem is Theorem~\ref{thm:intrinsic-fibre}.
Fix retained dimension $k$ and a positive reference representation $Q$ with $e$ endpoints. For each $0\le e'\le e$, the set of positive $e'$-endpoint representations of the same labelled function $G_Q$ is semialgebraic and has an explicit finite equality test. For a candidate $Q'$, form its active data $K_J',S_J'$ and the reference data $K_I,S_I$. For each pair $I,J$, let $v_1,\ldots,v_r$ generate the polyhedral cone $K_I\cap K_J'$. Then equality of the two functions is equivalent to
\begin{equation}\label{eq:all-strata-test}
 v_a^{\mathsf T}(S_I-S_J')v_b=0
 \quad\text{for every }I,J\text{ and }1\le a,b\le r.
\end{equation}
The zero cone imposes no condition. The least $e'$ for which these conditions and positive definiteness are feasible is the minimal endpoint dimension. The normalized minimal fibre is obtained by additionally setting $\diag C'=\mathbf1$; all unnormalized representations follow by positive endpoint scales, with permutations left as a finite ambiguity.

As the reference matrix varies in a semialgebraic family with fixed $k,e$, there is a finite semialgebraic partition on which the minimal endpoint dimension and the semialgebraic homeomorphism type of each normalized candidate fibre are constant. This includes arbitrary losses of visibility, changes of cross-block rank, and coalescence of active-set matrices.
\end{theorem}
\begin{proof}
The active formulas remain valid without full rank or strict visibility, because each principal endpoint block is positive definite. On $K_I\cap K_J'$ the two values are $x^{\mathsf T}S_Ix$ and $x^{\mathsf T}S_J'x$. These intersection cones cover $X$. If their difference vanishes on a cone, it vanishes at every generator and at every sum of two generators. Polarization gives all equalities in \eqref{eq:all-strata-test}. Conversely, expanding a nonnegative linear combination of the generators shows that those equalities make the quadratic difference vanish on the whole cone. This proves the exact finite test, including lower-dimensional intersections.

For semialgebraicity one may write the same test without choosing generators:
\[
 \bigwedge_{I,J}\ \forall x\;
 \bigl(x\in K_I\cap K_J'\Longrightarrow
 x^{\mathsf T}(S_I-S_J')x=0\bigr).
\]
Every coefficient is rational in the entries of $Q,Q'$, with denominators products of positive principal minors. Clearing these positive denominators and applying real quantifier elimination \cite{BCR} gives a finite Boolean combination of polynomial equalities and inequalities. Alternatively, extreme rays of each pointed intersection cone are obtained from finitely many choices of active linear equalities; partition by the relevant minors and their signs to obtain a finite generator implementation. No generic-rank assumption is used.

There is always a feasible candidate at $e'=e$, namely $Q$. Testing the finitely many smaller dimensions therefore determines the minimum. Positive diagonal scaling normalizes every positive endpoint block, so it loses no representation. Now regard each normalized equality set as a semialgebraic set over the reference parameter space. Semialgebraic triviality \cite{Hardt} applied to each of these finitely many projections gives finite partitions with constant fibre homeomorphism type, including the empty fibre. A common refinement fixes the least nonempty dimension and proves the final assertion.
\end{proof}

The last theorem is an effective presentation and a finite-stratification result, not an assertion that every higher-dimensional fibre is a product of the scalar intervals of Theorem~\ref{thm:three-cells}. Its semialgebraic triviality step is classical. The value-function equality test specifies the inverse object to which that theorem is applied; the first-stratum theorems supply explicit reconstruction where a much sharper description is possible.
